\section{Введение}

Краткая саммари того о чем будет говориться дальше - часто читают только ее. Объем не менее 5 страниц, включая введение в предметную область. Во введении обычно должны использоваться все термины из глоссария.

\subsection{Цель}
Цель работы заключается в - Ключевые аспекты экономик стран обеспечивающих успех бизнесу в сфере финансовых технологий

\subsection{Объект исследования} 
Экономические аспекты государств, на примере набора дынных из~\cite{}.

\subsection{Предмет исследования}
Экономические аспекты исследуются на предмет важности для успешного и устойчивого развития компаний в области финансовых технологий, по материалам набора данных из~\cite{?}.


\subsection{Методы исследования}
\begin{itemize}
    \item Программирование на Python с использованием библиотек:  \cite{?}, ....
    \item Визуализация данных с использованием Yandex DataLens.
    \item Обзор академических источников.
\end{itemize}

\subsection{Задачи}

%1) Найти рейтинг стран по FinTech команиям - по количеству по их объямам что-то - это дополнительный наборных
%2) У нас есть страны уже по рейтинка экномик-фридом.
%3) Соединить 1 и 2 - и получить ответ какие факторы в рейтиге экономик-фридом коррелируют с успешностью финтах компании. Тут гипотеза - что чем больше в стране финтеха тем они там успешнее.

Задачи проекта разделяются на две части: исследовательскую и программную, а именно:

\begin{itemize}
  \item{Исследовательская часть}
    \begin{itemize}
        \item {Исследовать дополнительные наборы данных по фин-тех компаниям}
        \item {Провести описательный анализ набора данных из \cite{?}}
        \item {Провести обзор источников, использующих рассматриваемый набор данных.}
        \item {Изучить возможность подключения дополнительного набора данных, который позволит выявлять страну популярности.}
        %...
    \end{itemize}
  \item{Программная часть}
      \begin{itemize}
        \item{Подготовить }
    \end{itemize}
\end{itemize}

\subsection{Актуальность работы}
Отдельно отметить актуальность работы и применение полученных результатов. Обычно заканчивается фразой: "Таким образом можно считать, что \textbf{актуальность работы} продемонстрирована".

\subsection{План работы}
% План из 6/10 задач - табличка из трех колонок - краткое название задачи - дата начала и дата окончания

% Табличка из трехколонок - краткое название задачи, дата начала и дата окончания.
